\documentclass[../../../include/open-logic-section]{subfiles}

\begin{document}

\olfileid{sth}{choice}{tarskiscott}
\olsection{टार्स्की--स्कॉटची युक्ती}

\olref[cardinals][cardsasords]{defcardinalasordinal} मध्ये संचांकांना क्रमसूचक संख्या म्हणून परिभाषित केले. यासाठी सुक्रमणाचे स्वयंसिद्धक गृहीत धरले होते. हीच प्रचलित प्रमाणपद्धत आहे, एवढ्याच कारणासाठी आपण तसे केले.

म्हणून सुक्रमणाशी संबंधित तात्त्विक प्रश्नांची चर्चा करण्याआधी, प्रचलित प्रमाणपद्धत सोडून सुक्रमण गृहीत \emph{न धरता} संचांकांची उपपत्ती विकसित करणे आपल्याला \emph{शक्य} आहे, हे स्पष्ट असणे महत्त्वाचे आहे. \olref[sfr][siz][]{chap} मध्ये आलेल्या $\cardeq{A}{B}$, $\cardle{A}{B}$ आणि $\cardless{A}{B}$ या व्याख्या तरीही वापरता येतील. फक्त \emph{संचांकाची} नवी संकल्पना लागेल.

$A$ चा संचांक पुढीलप्रमाणे परिभाषित करावा, अशी साधी कल्पना सुचू शकते:
\[
	\Setabs{x}{\cardeq{A}{x}}.
\]
\olref[cardinals][hp]{sec} च्या अगदी शेवटी मांडलेल्या फ्रेगेच्या $\# x Fx$ या व्याख्येशी याची तुलना करता येईल. तिथे सूचित केलेल्या कारणांमुळे ही सर्वसाधारण व्याख्या अपयशी ठरते. कोणताही एकघटकी संच $\{\emptyset\}$ शी तुल्यबल असतो. पण श्रेणीच्या प्रत्येक उत्तरवर्ती टप्प्यावर नवे एकघटकी संच तयार होतात (आधीच्या टप्प्याचा एकघटकी संचच पाहा). त्यामुळे अरिक्त आधारसंचाच्या बाबतीत $\Setabs{x}{\cardeq{A}{x}}$ अस्तित्वात नसतो, कारण त्याला प्रतांक असू शकत नाही. रिक्त आधारसंचाचे प्रकरण या आक्षेपाच्या कक्षेत येत नाही.

ही अडचण टाळण्यासाठी टार्स्की आणि स्कॉट यांची युक्ती वापरू:\footnote{आठवण: स्पष्टपणे वेगळे सांगितले नसेल, तर सर्व सूत्रांत परिमिती असू शकतात.}

\begin{defn}[टार्स्की--स्कॉट]
कोणत्याही सूत्र $\phi(x)$ साठी, $\phi(x)$ पूर्ण करणाऱ्या आणि शक्य तितक्या लघुतम प्रतांकाच्या सर्व $x$ चा संच $[ x : \phi(x)] $ असू द्या ($\phi$ पूर्ण करणारी कोणतीही वस्तू नसल्यास $\emptyset$ घ्या).
\end{defn}

ही व्याख्या योग्य आहे हे तपासले पाहिजे. $\ZF$ मध्ये काम करताना \olref[spine][foundation]{zfentailsregularity} नुसार प्रत्येक $x$ साठी $\setrank{x}$ अस्तित्वात असतो. आता $\phi$ पूर्ण करणाऱ्या काही वस्तू असतील, तर $(\exists x\subseteq V_\alpha)\phi(x)$ पूर्ण करणारा लघुतम प्रतांक $\alpha$ घेता येतो; त्याच्या लघुतमतेचा अर्थ $(\forall \beta
\in \alpha)(\forall x \subseteq V_\beta)\lnot \phi(x)$ असा आहे. त्यानंतर विभक्तीकरणाने $[x : \phi(x)]$ हा संच $\Setabs{x \in
V_{\alpha+1}}{\phi(x)}$ म्हणून परिभाषित करता येतो.

टार्स्की--स्कॉटच्या युक्तीचे समर्थन केले आहे; आता ती वापरून संचांकाची संकल्पना परिभाषित करता येईल:

\begin{defn}
$A$ चा \textsc{ts}-संचांक म्हणजे $\text{tsc}(A) = [x :
\cardeq{A}{x}]$.
\end{defn}

\textsc{ts}-संचांकाच्या व्याख्येत सुक्रमण वापरलेले नाही. पण ते स्वयंसिद्धक नसतानाही \emph{\textsc{ts}-संचांकांचे} वर्तन \olref[cardinals][cardsasords]{defcardinalasordinal} मध्ये परिभाषित केलेल्या \emph{संचांकांसारखेच} असल्याचे दाखवता येते. उदाहरणार्थ, \olref[cardinals][cardsasords]{lem:CardinalsBehaveRight} आणि \olref[card-arithmetic][opps]{lem:SizePowerset2Exp} हे निष्कर्ष \textsc{ts}-संचांकांच्या भाषेत पुन्हा मांडल्यास, सुक्रमण गृहीत न धरता त्यांचे पुरावे $\ZF$ मध्ये व्यवस्थित चालतात.

याच संदर्भात, टार्स्की--स्कॉटची युक्ती वापरून क्रमसूचक संख्यांची उपपत्तीही विकसित करता येते. $\tuple{A, <}$ हे सुक्रमण असताना $\text{tso}(A, <) = [\tuple{X,
R} : \ordeq{\tuple{A, <}}{\tuple{X, R}}]$ घ्या. संचांक आणि क्रमसूचक संख्या यांच्या या मांडणीविषयी अधिक माहितीसाठी \citet[chs.~9--12]{Potter2004} पाहा.

\end{document}
